\subsection{特征标与格罗滕迪克群}

用 \(\theta_{G}\) 表示从 \(K\otimes_{\mathbf{Z}}\mathrm{R}_{\mathrm{K}}(\mathrm{G})\) 到 \(\mathcal{Z}_{\mathrm{K}}(\mathrm{G})\) 的 K-代数同态，它把 \(1\otimes[\pi]\) 映为 \(\chi_{\pi}\)（对每个有限维表示 \(\pi\)）。

命题 8. —— a) 同态 \(\theta_{\mathrm{G}}\) 是从 \(\mathrm{K} \otimes_{\mathbf{Z}} \mathrm{R}_{\mathrm{K}}(\mathrm{G})\) 到 \(\mathcal{Z}_{\mathrm{K}}(\mathrm{G})\) 的同构。

\begin{enumerate}
\def\labelenumi{\alph{enumi})}
\setcounter{enumi}{1}
\tightlist
\item
  假设 K 的特征为 0。那么 \(\theta_{\mathrm{G}}\) 定义了从 \(\mathrm{R}_{\mathrm{K}}(\mathrm{G})\) 到 \(\mathcal{Z}_{\mathrm{K}}(\mathrm{G})\) 的一个子环的同构，该子环由诸特征标 \(\chi_{\lambda}\)（\(\lambda\) 取遍 \(\widehat{\mathrm{G}}\)）的整系数线性组合组成。
\end{enumerate}

族 \(([\lambda])_{\lambda \in \widehat{\mathbf{G}}}\) 是 \(\mathbf{Z}\)-模 \(\mathrm{R}_{\mathrm{K}}(\mathrm{G})\) 的一组基，而族 \((\chi_{\lambda})_{\lambda \in \widehat{\mathbf{G}}}\) 是 K-向量空间 \(\mathcal{Z}_{\mathrm{K}}(\mathrm{G})\) 的一组基。断言 a) 与 b) 由此得出（VIII, 第 411 页，命题 5）。

